\subsection{6.5 Improper Integrals}
\hfill 407
\medskip
c) $\text{Ci}x = \int_{x}^{+\infty} \frac{\cos t}{t} dt$ (the cosine integral) can be computed for sufficiently
large values of $x$ by the approximate formula $\text{Ci}x \approx \frac{\sin x}{x}$; estimate the region of
values where the absolute error of this approximation is less than $10^{-4}$.

\textbf{2.} Show that
$$
\int_{1}^{+\infty} \frac{\sin x}{x^\alpha} dx, \quad \int_{1}^{+\infty} \frac{\cos x}{x^\alpha} dx
$$
absolutely only for $\alpha > 1$;
b) the Fresnel integrals
$$
C(x)= \frac{1}{\sqrt{2}} \int_{0}^{x} \cos t^2 dt, \quad S(x)= \frac{1}{\sqrt{2}} \int_{0}^{x} \sin t^2 dt
$$
are infinitely differentiable functions on the interval $]0, +\infty[$, and both have a limit
as $x \to +\infty$.

\textbf{3.} Show that
a) the elliptic integral of first kind
$$
F(k, \varphi) := \int_{0}^{\varphi} \frac{dt}{\sqrt{(1-t^2)(1-k^2t^2)}}
$$
is defined for $0 \le k < 1, 0 \le \varphi \le \frac{\pi}{2}$ and can be brought into the form
$$
F(k, \varphi) = \int_{0}^{\varphi} \frac{d\psi}{\sqrt{1-k^2 \sin^2 \psi}};
$$
b) the complete elliptic integral of first kind
$$
K(k) = \int_{0}^{\pi/2} \frac{d\psi}{\sqrt{1-k^2 \sin^2 \psi}}
$$
increases without bound as $k \to 1-0$.

\textbf{4.} Show that
a) the exponential integral $\text{Ei}(x) = \int_{x}^{+\infty} \frac{e^t}{t} dt$ is defined and infinitely differen-
tiable for $x < 0$;
b) $-\text{Ei}(-x) = \frac{e^{-x}}{x} \left(1 - \frac{1}{x} + \frac{2!}{x^2} - \frac{3!}{x^3} + \dots + \frac{(-1)^n n!}{x^n} + o\left(\frac{1}{x^n}\right)\right)$ as $x \to +\infty$;
c) the series $\sum_{n=0}^{\infty} \frac{(-1)^n x^n}{n \cdot n!}$ does not converge for any value of $x \in \mathbb{R}$;
d) $\text{li}(x) \sim \frac{x}{\ln x}$ as $x \to +0$. (For the definition of the logarithmic integral $\text{li}(x)$
see Example 21.)

\textbf{5.} Show that
a) the function $\Phi(x) = \frac{1}{\sqrt{\pi}} \int_{0}^{x} e^{-t^2} dt$, called the error function and often de-
noted $\text{erf}(x)$, is defined, even, and infinitely differentiable on $\mathbb{R}$ and has a limit as
$x \to +\infty$;